(1) 在非中心区域 \(\mu y_{tt}=Ty_{xx}\)，波速 \(c=\sqrt{T/\mu}\)。弦两端 \(y(\pm L,t)=0\)，中心位移连续。点质量沿横向受两侧张力合力，满足 \(\boxed{M y_{tt}(0,t)=T[y_x(0^+,t)-y_x(0^-,t)]}\)，不可将点质量简单等同附加线密度而遗漏导数跳变。

(2) 令 \(\omega=ck\)。偶模在两侧可写 \(Y=A\sin[k(L-|x|)]\)，中心两侧导数为 \(\mp Ak\cos kL\)，代入跳跃条件得 \(\boxed{2T k\cot(kL)=M\omega^2}\)，条件在 \(\sin kL\ne0\) 下有效。奇模中心位移为零，满足 \(\sin kL=0\)，即 \(k_n=n\pi/L\)，完全不加载点质量动能。

(3) 若 \(M\to0\)，偶模条件趋 \(\cot kL=0\)，即 \(kL=(n+1/2)\pi\)，与奇模合并成总长 \(2L\) 的完整固定弦频谱。若 \(M\to\infty\)，正频偶模向 \(\sin kL=0\) 靠近；中心重物近似固定节点，出现两半弦的谱重合，另有很低的质量主导模。

(4) 全体系守恒能量 \(E=\frac12\int_{-L}^L[\mu y_t^2+Ty_x^2]dx+\frac M2y_t(0)^2\)。乘两个本征方程交叉积分并用中心力平衡与固定端条件，得到若 \(\omega_n^2\ne\omega_m^2\)，则 \(\int_{-L}^L\mu Y_nY_mdx+M Y_n(0)Y_m(0)=0\)。因此正交权重中必须包括点质量的Delta分布。
